% common/user-install.tex
% Shared chapter of the user part: installing and compiling.
% Uses \uniname, \unishort, \uniid, set by the guide.
\chapter{Installing and compiling}
\label{ch:install}

\section{Requirements}

\begin{itemize}
  \item A recent \TeX{} distribution: \TeX{}~Live or Mac\TeX{} 2023 or
    later, or an up-to-date MiK\TeX. The class is written in \LaTeX3 and
    needs a modern \LaTeX{} kernel.
  \item \textbf{pdf\LaTeX} is the reference engine, used by all the
    examples in this guide. Xe\LaTeX{} and Lua\LaTeX{} are also accepted
    (the class then loads \pkg{fontspec}).
  \item For the bibliography, \pkg{biblatex} and \pkg{biber}, both part of
    every full \TeX{} distribution. Without them the class falls back to
    \pkg{natbib} and Bib\TeX{} (Section~\ref{sec:bib}).
\end{itemize}

\section{The files you need}

\pkg{academicthesis} is a collection of small packages that load one
another \emph{by name}. For a \unishort{} thesis you need only these:

\begin{tabularx}{\textwidth}{@{}l Y@{}}
\toprule
\textbf{Files} & \textbf{Role} \\
\midrule
\file{core/academicthesis.cls}, \file{core/academicthesis-kernel.sty} & The core class, shared by every university. \\
\file{languages/academicthesis-<language>.ldf} & The words printed on the pages, one file per language (English, Italian, French). \\
\texttt{universities/\uniid/*} & The \unishort{} package: its layout, labels and logos. \\
\bottomrule
\end{tabularx}

The other folders of \file{universities/} belong to other universities and
can be left out.

\section{Three ways to set it up}

\begin{enumerate}
  \item \textbf{Inside the repository.} Copy the \unishort{} showcase in
    \file{showcases/} and keep its first line,
    \begin{lstlisting}
\input{academicthesis-repo-paths}
    \end{lstlisting}
    which tells \LaTeX{} where the package folders are. Then compile with
    any editor, or from a terminal with \texttt{latexmk}.
  \item \textbf{Next to your thesis.} Copy the files listed above into the
    folder of your thesis (without sub-folders) and delete the
    \verb|\input{academicthesis-repo-paths}| line. This is also how to
    use the class on Overleaf: upload the files to the project.
  \item \textbf{Installed in your personal \TeX{} tree}, for example
    \file{~/Library/texmf/tex/latex/academicthesis/} on macOS or
    \file{~/texmf/tex/latex/academicthesis/} on Linux. Every document on the
    computer can then use the class, with no extra line.
\end{enumerate}

\begin{notebox}{Editors}
Texifier/Texpad, TeXShop, TeXstudio, VS Code (LaTeX Workshop) and Overleaf
all work. In Texifier/Texpad, set the typesetter to \emph{TeX Live
(pdfLaTeX)}: its built-in engine is too old for this class.
\end{notebox}

\section{Compiling}

Run pdf\LaTeX{} twice (for the table of contents and cross-references),
with \pkg{biber} between the two runs when the thesis has a bibliography.
\texttt{latexmk} does all of this by itself:

\begin{lstlisting}[style=shell]
latexmk -pdf my-thesis.tex
\end{lstlisting}
